\section{Literature Review}

McMahan et al. [1] introduced the Federated Averaging (FedAvg) algorithm at AISTATS 2017, achieving 10-100$\times$ reduction in communication rounds compared to federated SGD with near-centralized accuracy on MNIST and CIFAR-10. However, FedAvg suffers from weight divergence with highly skewed non-IID data, experiencing significant performance degradation with extreme data heterogeneity common in healthcare scenarios. Building upon this, Wang et al. [2] proposed FedMA at ICLR 2020, using the Hungarian algorithm for neuron matching before averaging, achieving 5-15\% accuracy improvement over FedAvg on heterogeneous datasets. Despite these gains, FedMA introduces 30-50\% additional training time and remains limited to CNNs and LSTMs, struggling with transformer architectures.

Kairouz et al. [3] provided a comprehensive survey in Foundations and Trends in Machine Learning 2021, identifying key challenges including data heterogeneity, privacy leakage, fairness, and system heterogeneity while analyzing 150+ federated methods. Though comprehensive, this survey focused on general federated learning without addressing healthcare-specific requirements or production deployment challenges. Liu et al. [4] addressed production concerns with FATE, the first industrial-grade federated platform supporting Private Set Intersection, homomorphic encryption, and cross-cloud deployment in JMLR 2021. However, FATE's complex Kubernetes-based setup and lack of healthcare-specific features create barriers for medical institutions.

Beutel et al. [5] introduced Flower, a lightweight and extensible framework supporting multiple ML backends with modular design for custom aggregation strategies. While excellent for research, Flower lacks healthcare compliance features, automated audit logging, and production-ready dashboards necessary for clinical deployment. Sheller et al. [6] demonstrated federated learning's medical viability in Scientific Reports 2020, achieving 99\% centralized accuracy on brain tumor segmentation across 10 institutions. Despite proving feasibility, their approach required manual coordination and significant technical expertise at each site, limiting scalability.

Haripriya et al. [7] proposed adaptive aggregation combining FedAvg and FedSGD strengths in Scientific Reports 2025, reducing training rounds by 15-20\% through dynamic adjustment based on convergence metrics. However, this approach introduces 20-25\% computational overhead for larger models with limited real-world deployment evaluation. Wei et al. [8] provided comprehensive differential privacy analysis in IEEE TIFS 2020, demonstrating privacy-utility tradeoffs with 3-8\% accuracy degradation using moment accountant methods. Determining appropriate privacy budgets for healthcare remains challenging, and noise addition significantly impacts convergence.

Li et al. [9] introduced FedProx with a proximal term enabling partial work from resource-constrained clients, maintaining convergence even when clients complete only 10-50\% of assigned work. While addressing system heterogeneity, FedProx requires additional hyperparameter tuning and focuses less on data heterogeneity challenges. Mo et al. [10] leveraged Intel SGX for hardware-based security in MobiSys 2021, achieving stronger guarantees than cryptographic methods with only 5-10\% overhead compared to 200-300\% for homomorphic encryption. However, TEE-based approaches depend on specific hardware availability and face known vulnerabilities from side-channel attacks.

\textbf{Summary of Findings:} The literature reveals substantial theoretical progress but critical gaps for healthcare adoption. FedAvg and FedMA struggle with non-IID medical data, while existing frameworks (Flower, FATE) lack healthcare-ready features like automated compliance and intuitive interfaces. Privacy mechanisms introduce significant tradeoffs—differential privacy reduces accuracy by 3-8\% while homomorphic encryption adds 200-300\% overhead. Medical collaborations prove federated learning's feasibility (99\% accuracy) but require manual coordination. No current platform combines automated HIPAA/GDPR compliance, real-time monitoring for non-technical users, healthcare-optimized algorithms, and simplified deployment.

\section{Functional Requirements}

\textbf{Core Training Features (FR1-FR4):} The platform coordinates distributed training through a central coordinator that broadcasts initial models, manages training rounds, collects encrypted updates, performs weighted aggregation, and distributes improved models. Local hospitals train ResNet-50 models on private medical imaging datasets using PyTorch with data augmentation, computing model updates without transmitting raw data. All communications use HTTPS/TLS encrypted channels and WebSocket protocols with authentication and integrity checks. The system supports multiple aggregation algorithms including FedAvg (baseline), FedOpt (adaptive learning rates), FedMA (neuron alignment), FedProx (system heterogeneity), and a novel algorithm optimized for non-IID medical data.

\textbf{Monitoring and Privacy (FR5-FR7):} Real-time monitoring dashboards provide live visualization of training progress, model metrics, and system health via WebSocket updates, displaying round numbers, accuracy/loss trends, client status, and convergence indicators. The platform guarantees raw patient data never leaves local servers, transmitting only model weights with automated privacy checks and comprehensive audit logs. Medical image preprocessing follows standardized pipelines for resizing (224$\times$224), normalization, metadata removal, data augmentation, and quality validation.

\textbf{Advanced Management (FR8-FR11):} Dynamic hospital management supports node addition/removal during training with graceful disconnection handling and mid-training onboarding. Model versioning maintains complete history with automatic checkpointing after each round, supporting rollback, export, and lineage tracking. Performance benchmarking tracks accuracy, convergence speed, communication overhead, and training time across algorithms. Compliance audit logging maintains tamper-evident records of all connection events, model transmissions, aggregation operations, and configuration changes for HIPAA/GDPR audits.

\textbf{Access Control and Operations (FR12-FR15):} Role-based access control supports hospital administrators, data scientists, viewers, and coordinator administrators with institutional identity provider integration. Error handling includes automatic retry for network failures, timeout detection for stragglers, model update validation, critical error alerting, and degraded operation modes. Configuration management provides flexible settings for training hyperparameters, aggregation algorithms, client sampling, communication protocols, and preprocessing parameters. Export capabilities generate trained model weights (PyTorch, ONNX), training metrics (CSV), visualization plots, compliance reports, and configuration snapshots.

\section{Hardware and Software Requirements}

\subsection{Hardware Requirements}

The platform is designed to accommodate a range of institutional resources, from smaller clinics with standard workstations to large hospitals with dedicated GPU servers. The following table outlines the minimum and recommended specifications.

\begin{table}[H]
\centering
\caption{Hardware Specifications}
\begin{tabularx}{\textwidth}{|l|X|X|}
\hline
\textbf{Component} & \textbf{Coordinator Server} & \textbf{Hospital Client} \\ \hline
Processor & 8+ core CPU (16+ recommended) & 4+ core CPU (Intel i7/AMD R7 rec.) \\ \hline
Memory & 16GB RAM (32GB recommended) & 8GB RAM (16GB recommended) \\ \hline
Storage & 100GB SSD (500GB recommended) & 50GB SSD storage \\ \hline
Network & 100 Mbps (1 Gbps recommended) & 10 Mbps upload (50+ Mbps rec.) \\ \hline
GPU & N/A & Optional (NVIDIA 6GB+ VRAM rec.) \\ \hline
\end{tabularx}
\end{table}

\subsection{Software Requirements}

The development and deployment of the platform rely on a modern, containerized software stack that ensures consistency across different hospital environments.

\begin{table}[H]
\centering
\caption{Software Stack}
\begin{tabularx}{\textwidth}{|l|X|}
\hline
\textbf{Category} & \textbf{Software / Framework} \\ \hline
Core Frameworks & Flower v1.5+, PyTorch v2.0+, FastAPI v0.104+, Next.js 15, Docker v24.0+ \\ \hline
Languages & Python 3.9+, TypeScript 5.0+, JavaScript ES2020+ \\ \hline
Database & Supabase (PostgreSQL 14+) \\ \hline
Development & Node.js 18 LTS+, Git, npm/yarn, pip \\ \hline
Libraries & NumPy 1.24+, Pillow 10.0+, torchvision 0.15+, pyjwt, bcrypt \\ \hline
GPU Support & NVIDIA CUDA Toolkit 11.8+, cuDNN 8.6+ \\ \hline
\end{tabularx}
\end{table}

\section{Summary}

The literature review reveals federated learning's theoretical maturity but identifies critical gaps for healthcare deployment. While algorithms like FedAvg provide communication efficiency and FedMA improves aggregation quality, both struggle with non-IID medical data. Existing frameworks lack healthcare-specific features including automated compliance, intuitive interfaces, and simplified deployment. Medical collaboration studies validate federated learning's feasibility (99\% centralized accuracy) but rely on manual coordination requiring extensive technical expertise.

Our platform addresses these gaps through comprehensive orchestration supporting multiple aggregation algorithms, real-time WebSocket dashboards, automated compliance logging, role-based access control, and dynamic hospital management. Hardware requirements accommodate resource-constrained hospitals (minimum 8GB RAM, 4-core CPU) while supporting GPU acceleration for well-equipped institutions, democratizing federated learning access. The software stack—Flower, PyTorch, FastAPI, Next.js, Docker—provides production-grade reliability and deployment flexibility, positioning this work to deliver the first healthcare-optimized federated learning platform bridging the gap between research demonstrations and clinical deployment.

